\documentclass[11pt,a4paper]{article}
\usepackage{amsmath,amssymb,amsthm}
\usepackage{authblk}
\usepackage[margin=2.5cm]{geometry}
\usepackage{hyperref}
\usepackage{booktabs}
\usepackage{url}
\usepackage{enumerate}

%% Theorem environments
\newtheorem{theorem}{Theorem}[section]
\newtheorem{proposition}[theorem]{Proposition}
\newtheorem{corollary}[theorem]{Corollary}
\newtheorem{lemma}[theorem]{Lemma}
\newtheorem{conjecture}[theorem]{Conjecture}
\theoremstyle{definition}
\newtheorem{definition}[theorem]{Definition}
\newtheorem{remark}[theorem]{Remark}
\newtheorem{example}[theorem]{Example}

%% Macros
\newcommand{\ZZ}{\mathbb{Z}}
\newcommand{\QQ}{\mathbb{Q}}
\newcommand{\FF}{\mathbb{F}}
\newcommand{\CC}{\mathbb{C}}
\newcommand{\MPMNS}{M_{\mathrm{PMNS}}}
\newcommand{\MCKM}{M_{\mathrm{CKM}}}
\newcommand{\vol}{\mathrm{vol}}
\newcommand{\ITF}{\mathrm{ITF}}
\newcommand{\pk}{p_k}
\newcommand{\Xo}{X_0(11)}
\newcommand{\Ek}{E(\FF_{\pk})}

\begin{document}

\title{The $\ZZ/5$ Bridge:\\
Hyperbolic Covering Towers, Rational Torsion of $X_0(11)$,\\
and an Arithmetic Proof of the Universal Cover Prime Formula}

\author{Marvin L. Gentry}
\affil{Independent Researcher, Seattle, WA, USA\\
\texttt{drmlgentry@protonmail.com}\quad
ORCID: 0009-0006-4550-2663}
\date{\today}
\maketitle

%%------------------------------------------------------------------
\begin{abstract}
%%------------------------------------------------------------------
We establish a closed chain of implications connecting three
independent mathematical structures:
the degree-5 covering towers of two compact arithmetic hyperbolic
3-manifolds, the rational 5-torsion of the modular curve $X_0(11)$,
and the Hecke eigenvalue sequence of the associated newform.

Let $M_k^{\mathrm{PMNS}} = \texttt{m003}(-(k+1), 2k+1)$ and
$M_k^{\mathrm{CKM}} = \texttt{m006}(1-3(k+1), k+1)$ be the
$k$-th elements of the PMNS and CKM Farey towers, each with
$H_1 \cong \ZZ/5$.
Via Reidemeister--Schreier theory we prove that for $k = 1, \ldots, 7$,
the unique degree-5 cover of each tower element has
$H_1 \cong (\ZZ/\pk)^2$ where $\pk = 5k(k+1)+1$.

The central observation is that $\pk \equiv 1 \pmod{5}$
for all $k$, a trivial algebraic identity.
Combined with the fact that $X_0(11)$ has rational 5-torsion
$X_0(11)(\QQ)[5] \cong \ZZ/5$, this forces:
\begin{enumerate}[(i)]
\item $\Ek \supseteq (\ZZ/5)^2$ for every prime $\pk$ with $k \geq 2$;
\item the Hecke eigenvalue satisfies $a_{\pk} \equiv 2 \pmod{5}$,
verified for all prime $\pk$ through $k=15$;
\item a canonical injection $\ZZ/5 \hookrightarrow \ZZ/\pk$
through which the covering homology and the elliptic curve
torsion are arithmetically linked.
\end{enumerate}
The $\ZZ/5$ torsion is the common arithmetic language between the
hyperbolic topology of the flavor manifolds and the modular arithmetic
of $X_0(11)$, whose conductor 11 equals the first cover prime $p_1$.
\end{abstract}

%%------------------------------------------------------------------
\section{Introduction}
\label{sec:intro}
%%------------------------------------------------------------------

The Hyperbolic Flavor Geometry (HFG) programme
\cite{GentrySSRN,GentryEM,GentryFarey}
identifies two compact arithmetic hyperbolic 3-manifolds
as geometric models for the Standard Model flavor parameters:
the Meyerhoff manifold $\MPMNS = \texttt{m003}(-2,3)$
for the PMNS lepton mixing matrix and
$\MCKM = \texttt{m006}(-5,2)$ for the CKM quark mixing matrix.
Both satisfy $H_1 \cong \ZZ/5$ and sit at the first position
of infinite Farey towers of Dehn fillings.

The present paper identifies a precise arithmetic bridge
linking these manifolds to the modular curve $\Xo$.
The bridge is the $\ZZ/5$ torsion --- a structure that appears
simultaneously in three places:
\begin{enumerate}[(a)]
\item as the first homology $H_1 \cong \ZZ/5$ of every tower element
      $M_k^{\mathrm{PMNS}}$ and $M_k^{\mathrm{CKM}}$;
\item as the rational torsion $\Xo(\QQ)[5] \cong \ZZ/5$
      generated by the cusp at $0$; and
\item as the common factor linking the cover prime sequence
      $\pk = 5k(k+1)+1 \equiv 1 \pmod{5}$ to the Hecke
      eigenvalue congruence $a_{\pk} \equiv 2 \pmod{5}$.
\end{enumerate}

\begin{theorem}[Main --- Closed Chain]
\label{thm:main}
For the cover prime formula $\pk = 5k(k+1)+1$, the following
chain of implications holds:
\[
  \pk \equiv 1 \pmod{5}
  \;\Rightarrow\;
  \Ek \supseteq (\ZZ/5)^2
  \;\Rightarrow\;
  a_{\pk} \equiv 2 \pmod{5}
  \;\Rightarrow\;
  \ZZ/5 \hookrightarrow H_1(M_k^{(5)}).
\]
Each implication is a proved theorem or elementary identity.
\end{theorem}

The paper is organised as follows.
Section~\ref{sec:setup} recalls the two Farey towers and proves
the cover prime formula for both via Reidemeister--Schreier theory.
Section~\ref{sec:x011} establishes the torsion structure of $\Xo$
and its reduction at the cover primes.
Section~\ref{sec:chain} assembles the closed chain and proves
Theorem~\ref{thm:main}.
Section~\ref{sec:identity} establishes the algebraic identity
$\pk = L_k^2 + k$ and the forbidden Lucas prime theorem.
Section~\ref{sec:conjectures} collects open problems.

%%------------------------------------------------------------------
\section{The Two Farey Towers and Cover Primes}
\label{sec:setup}
%%------------------------------------------------------------------

\subsection{Surgery formulas}

\begin{theorem}[Surgery formulas]
\label{thm:surgery}
For all coprime $(p,q)$ with hyperbolic Dehn fillings:
\begin{align}
  |H_1(\texttt{m003}(p,q))| &= 5|2p+q|, \label{eq:surgery-pmns}\\
  |H_1(\texttt{m006}(p,q))| &= 5|p+3q|. \label{eq:surgery-ckm}
\end{align}
\end{theorem}

\begin{proof}
The fundamental group of $\texttt{m003}$ has relator $R = \texttt{aabAbAB}$.
Abelianization gives $H_1(\texttt{m003}) \cong \ZZ/5 \oplus \ZZ$.
The peripheral curves abelianize to $[\mu] = 2[a] - [b]$
and $[\lambda] = [a] - 2[b]$.
Dehn surgery along $(p,q)$ kills $p[\mu]+q[\lambda]$,
yielding determinant $5(2p+q)$.

For $\texttt{m006}$, the relator is $\texttt{ababbAAbb}$.
Abelianization gives $0 \cdot [a] + 5[b] = 0$,
so $H_1(\texttt{m006}) \cong \ZZ \oplus \ZZ/5$.
The peripheral curves abelianize to
$[\mu] = -[a]-2[b]$ and $[\lambda] = -3[a]+[b]$,
giving determinant $5(p+3q)$.
\end{proof}

\subsection{The two Farey towers}

\begin{definition}
The \emph{PMNS Farey tower} is
$M_k^{\mathrm{PMNS}} = \texttt{m003}(-(k+1), 2k+1)$ for $k \geq 1$.
The \emph{CKM Farey tower} is
$M_k^{\mathrm{CKM}} = \texttt{m006}(1-3(k+1), k+1)$ for $k \geq 1$.
\end{definition}

\begin{proposition}
\label{prop:h1}
$H_1(M_k^{\mathrm{PMNS}}) \cong H_1(M_k^{\mathrm{CKM}}) \cong \ZZ/5$
for all $k \geq 1$. Both towers have Farey determinant $1$
between consecutive slope pairs.
\end{proposition}

\begin{proof}
For PMNS: $|2(-(k+1))+(2k+1)| = 1$,
so $|H_1| = 5$ by~\eqref{eq:surgery-pmns}.
For CKM: $|(1-3(k+1))+3(k+1)| = 1$,
so $|H_1| = 5$ by~\eqref{eq:surgery-ckm}.
The Farey determinant calculation is direct.
\end{proof}

\subsection{Cover primes via Reidemeister--Schreier}

Each manifold $M_k$ in either tower has a SnapPy
two-generator presentation
$\pi_1(M_k) = \langle a, b \mid r_1, r_2(k) \rangle$.

\begin{lemma}[H1 map]
\label{lem:phi}
The map $\varphi\colon \pi_1(M_k) \to \ZZ/5$ defined by
$\varphi(a) = 1$, $\varphi(b) = 2$ is the unique surjection
for every element of both towers.
\end{lemma}

Let $S = \ker\varphi$ be the unique index-5 subgroup.
Coset representatives are $\{1, b, b^2, b^3, b^4\}$.
The Reidemeister--Schreier generators are
$A_i = b^i a b^{-i}$ for $i = 0,1,2,3,4$
(all $b$-type generators are trivial since $b$ acts by coset shift).
For each relator and each starting coset, rewriting produces
a linear relation in $H_1(S; \ZZ)$.
The full relation matrix $\mathbf{R}_k$ is $10 \times 5$.

\begin{theorem}[Cover homology formula]
\label{thm:cover}
For $k = 1, \ldots, 7$ and both towers,
$\mathbf{R}_k$ has Smith Normal Form diagonal
$[1, 1, 1, \pk, \pk]$, hence
\[
  H_1(S; \ZZ) \cong (\ZZ/\pk)^2,
  \qquad \pk = 5k(k+1)+1.
\]
\end{theorem}

\begin{proof}
Direct computation in SageMath for each $k$.
For $k=1$ and the PMNS tower, the relators are
$r_1 = \texttt{ababAbbAb}$ and
$r_2(1) = \texttt{abAbaabAbaBAB}$.
The $10 \times 5$ relation matrix has Smith Normal Form
$\mathrm{diag}(1,1,1,11,11)$, giving $H_1 \cong (\ZZ/11)^2$.
The computations for $k=2,\ldots,7$ and the CKM tower
proceed analogously; the results match in every case.
See Table~\ref{tab:cover}.
\end{proof}

\begin{table}[h]
\centering
\caption{Cover prime verification for both towers, $k=1,\ldots,7$.}
\label{tab:cover}
\begin{tabular}{rrccc}
\toprule
$k$ & $\pk$ & PMNS manifold & CKM manifold & Cover $H_1$ \\
\midrule
1 & 11  & $\texttt{m003}(-2,3)$  & $\texttt{m006}(-5,2)$  & $(\ZZ/11)^2$  \\
2 & 31  & $\texttt{m003}(-3,5)$  & $\texttt{m006}(-8,3)$  & $(\ZZ/31)^2$  \\
3 & 61  & $\texttt{m003}(-4,7)$  & $\texttt{m006}(-11,4)$ & $(\ZZ/61)^2$  \\
4 & 101 & $\texttt{m003}(-5,9)$  & $\texttt{m006}(-14,5)$ & $(\ZZ/101)^2$ \\
5 & 151 & $\texttt{m003}(-6,11)$ & $\texttt{m006}(-17,6)$ & $(\ZZ/151)^2$ \\
6 & 211 & $\texttt{m003}(-7,13)$ & $\texttt{m006}(-20,7)$ & $(\ZZ/211)^2$ \\
7 & 281 & $\texttt{m003}(-8,15)$ & $\texttt{m006}(-23,8)$ & $(\ZZ/281)^2$ \\
\bottomrule
\end{tabular}
\end{table}

\begin{remark}
The second relator of the PMNS tower has a recursive structure:
$r_2(k) = C^{k-1} \cdot B$ where
$C = \texttt{abAbaBAB}$ (a fixed 8-letter chunk) and
$B = \texttt{abAbaabAbaBAB}$ is a base word.
Each copy of $C$ contributes a closed loop in the covering
that increases the invariant factor by $p_{k+1} - p_k = 10k+6$.
This recursive structure is the combinatorial engine of
the cover prime formula.
\end{remark}

%%------------------------------------------------------------------
\section{The Modular Curve $X_0(11)$ and Its Reductions}
\label{sec:x011}
%%------------------------------------------------------------------

\subsection{Rational 5-torsion}

The modular curve $\Xo$ is an elliptic curve of conductor 11
with Weierstrass model
\[
  y^2 + y = x^3 - x^2 - 10x - 20.
\]
Its rational points form a group $\Xo(\QQ) \cong \ZZ/5$,
generated by the cusp at $\infty$.

\begin{proposition}
\label{prop:torsion}
$\Xo(\QQ)_{\mathrm{tors}} \cong \ZZ/5$,
generated by the cusp at $0$.
This torsion point has order exactly $5$ over $\QQ$.
\end{proposition}

This is classical; see for instance \cite[Example~8.1]{SilvermanAEC}.
The key consequence is that for any prime $p \neq 5, 11$,
the reduction $\bar{E}_p = \Xo \bmod p$ is an elliptic curve
over $\FF_p$ and the torsion point reduces to a point of order
(dividing) $5$ in $\bar{E}_p(\FF_p)$.

\subsection{Reduction at cover primes}

\begin{theorem}[5-torsion at cover primes]
\label{thm:torsion}
For every prime $\pk = 5k(k+1)+1$ with $k \geq 2$:
\begin{enumerate}[(i)]
\item $\pk \equiv 1 \pmod{5}$,
\item $5 \mid \#\Ek$,
\item $\Ek \supseteq (\ZZ/5)^2$,
\item $a_{\pk} \equiv 2 \pmod{5}$.
\end{enumerate}
\end{theorem}

\begin{proof}
(i) is the algebraic identity $\pk = 5k(k+1)+1 \equiv 1 \pmod{5}$.

For (ii): By Proposition~\ref{prop:torsion}, the rational 5-torsion
point of $\Xo$ reduces to a point of order 5 in $\bar{E}_{\pk}(\FF_{\pk})$
(since $\pk \neq 5, 11$). Hence $5 \mid \#\Ek$.

For (iii): Since $\pk \equiv 1 \pmod{5}$, the field $\FF_{\pk}$
contains the 5th roots of unity $\mu_5$.
An elliptic curve over $\FF_p$ with a rational 5-torsion point
and $p \equiv 1 \pmod{5}$ has full 5-torsion $E[5] \cong (\ZZ/5)^2$
defined over $\FF_p$, by the Weil pairing
(since $\mu_5 \subset \FF_p$).

For (iv): From (ii), $5 \mid (p+1-a_p)$.
Since $p \equiv 1 \pmod{5}$, we have $p+1 \equiv 2 \pmod{5}$.
Therefore $a_p \equiv 2 \pmod{5}$.
\end{proof}

\begin{table}[h]
\centering
\caption{Verification of Theorem~\ref{thm:torsion} for prime cover primes.}
\label{tab:hecke}
\begin{tabular}{rrcrrcc}
\toprule
$k$ & $\pk$ & prime & $a_{\pk}$ & $a_{\pk} \bmod 5$ & $\#\Ek$ & $(\ZZ/5)^2 \subset E$? \\
\midrule
1 & 11  & yes & 1   & 1 & 11  & no (bad reduction) \\
2 & 31  & yes & 7   & 2 & 25  & yes \\
3 & 61  & yes & 12  & 2 & 50  & yes \\
4 & 101 & yes & 2   & 2 & 100 & yes \\
5 & 151 & yes & 2   & 2 & 150 & yes \\
6 & 211 & yes & 12  & 2 & 200 & yes \\
7 & 281 & yes & $-18$ & 2 & 300 & yes \\
8 & 361 & no  & --- & --- & --- & (composite, skip) \\
11 & 661 & yes & 37 & 2 & 625 & yes \\
13 & 911 & yes & 12 & 2 & 900 & yes \\
14 & 1051& yes & 2  & 2 & 1050& yes \\
15 & 1201& yes & 2  & 2 & 1200& yes \\
\bottomrule
\end{tabular}
\end{table}

Note that $p_1 = 11$ is the conductor of $\Xo$ and is a point of
bad reduction; the theorem applies for $k \geq 2$.

%%------------------------------------------------------------------
\section{The Closed Chain}
\label{sec:chain}
%%------------------------------------------------------------------

We now assemble all pieces into the main theorem.

\begin{proof}[Proof of Theorem~\ref{thm:main}]
We verify each implication.

\textbf{Step 1:} $\pk \equiv 1 \pmod{5}$.
This is the algebraic identity of Theorem~\ref{thm:torsion}(i).

\textbf{Step 2:} $\pk \equiv 1 \pmod{5}
\Rightarrow \Ek \supseteq (\ZZ/5)^2$.
This is Theorem~\ref{thm:torsion}(iii), using the Weil pairing
and the rational torsion of $\Xo$.

\textbf{Step 3:} $\Ek \supseteq (\ZZ/5)^2
\Rightarrow a_{\pk} \equiv 2 \pmod{5}$.
From $5 \mid \#\Ek = \pk + 1 - a_{\pk}$ and $\pk \equiv 1 \pmod{5}$,
we get $a_{\pk} \equiv 2 \pmod{5}$.
This is Theorem~\ref{thm:torsion}(iv).

\textbf{Step 4:} $H_1(M_k^{(5)}) \cong (\ZZ/\pk)^2$.
This is Theorem~\ref{thm:cover}, proved by Reidemeister--Schreier.

\textbf{Step 5:} $\pk \equiv 1 \pmod{5}
\Rightarrow \ZZ/5 \hookrightarrow \ZZ/\pk$.
Since $\pk \equiv 1 \pmod{5}$, the group $\ZZ/\pk$ contains
a unique subgroup of order 5. This gives the canonical injection
$\ZZ/5 \hookrightarrow \ZZ/\pk$, and hence
$(\ZZ/5)^2 \hookrightarrow (\ZZ/\pk)^2 = H_1(M_k^{(5)})$.
\end{proof}

The closed chain can be displayed compactly as:
\begin{align*}
  \pk = 5k(k+1)+1
  &\Rightarrow \pk \equiv 1 \pmod{5} \\
  &\Rightarrow \FF_{\pk} \supseteq \mu_5,\quad
     \Xo(\QQ)[5] \cong \ZZ/5 \text{ reduces} \\
  &\Rightarrow \Ek \supseteq (\ZZ/5)^2 \\
  &\Rightarrow a_{\pk} \equiv 2 \pmod{5} \\
  &\Rightarrow \ZZ/5 \hookrightarrow H_1(M_k^{(5)}) = (\ZZ/\pk)^2.
\end{align*}

\begin{corollary}[Shared torsion structure]
\label{cor:shared}
The $\ZZ/5$ torsion acts simultaneously on the topological
covering $M_k^{(5)}$ (via $H_1$) and on the arithmetic
covering $\Xo[5](\FF_{\pk})$ (via the Weil pairing),
with a canonical isomorphism between the two $\ZZ/5$ factors
induced by the natural inclusion $\ZZ/5 \hookrightarrow \ZZ/\pk$.
\end{corollary}

%%------------------------------------------------------------------
\section{Algebraic Identities}
\label{sec:identity}
%%------------------------------------------------------------------

\subsection{Slope norm identity}

The PMNS slope at position $k$ is $(-(k+1), 2k+1)$,
with squared Euclidean norm
\[
  L_k^2 = (k+1)^2 + (2k+1)^2
         = 5k^2 + 4k + 1.
\]

\begin{theorem}[One-line identity]
\label{thm:identity}
$\pk = L_k^2 + k$ for all $k \geq 1$.
\end{theorem}

\begin{proof}
$\pk - L_k^2 = (5k^2+5k+1) - (5k^2+4k+1) = k$.
\end{proof}

\subsection{Forbidden Lucas primes}

\begin{theorem}[Forbidden Lucas primes]
\label{thm:forbidden}
The Lucas primes $L_2 = 2$, $L_3 = 3$, $L_4 = 7$
do not divide $\pk$ for any $k \geq 1$.
\end{theorem}

\begin{proof}
\textit{Mod 2:} $k(k+1)$ is always even,
so $5k(k+1) \equiv 0 \pmod{2}$ and $\pk \equiv 1 \pmod{2}$.

\textit{Mod 3:} $5 \equiv 2 \pmod{3}$ and
$k(k+1) \bmod 3 \in \{0, 2\}$ for $k \equiv 0, 1, 2 \pmod{3}$
respectively, giving $\pk \equiv 1$ or $2 \pmod{3}$.

\textit{Mod 7:} Direct computation of $\pk \bmod 7$
for $k = 0, \ldots, 6$ yields $\{1, 4, 3, 5, 3, 4, 1\}$,
none of which is $0$.
\end{proof}

\begin{theorem}[First permitted Lucas prime]
\label{thm:first-lucas}
The smallest Lucas prime dividing some $\pk$ is $L_5 = 11$,
which divides $\pk$ if and only if $k \equiv 1$ or $9 \pmod{11}$.
\end{theorem}

\begin{proof}
The congruence $5k(k+1)+1 \equiv 0 \pmod{11}$ reduces to
$k(k+1) \equiv 2 \pmod{11}$, solved by $k \equiv 1$ and
$k \equiv 9 \pmod{11}$.
At $k=1$: $p_1 = 11 = L_5$ itself.
At $k=9$: $p_9 = 451 = 11 \times 41$.
\end{proof}

\begin{corollary}
Both canonical manifolds $\MPMNS = M_1^{\mathrm{PMNS}}$ and
$\MCKM = M_1^{\mathrm{CKM}}$ share the first cover prime
$p_1 = 11 = L_5$, the conductor of $\Xo$.
\end{corollary}

%%------------------------------------------------------------------
\section{Universality}
\label{sec:universality}
%%------------------------------------------------------------------

\begin{theorem}[Universality criterion]
\label{thm:universality}
The cover prime formula $\pk = 5k(k+1)+1$ holds for any
cusped hyperbolic 3-manifold $M_{\mathrm{cusp}}$ with
$H_1(M_{\mathrm{cusp}}) \cong \ZZ/5 \oplus \ZZ$ and
a Farey tower of Dehn fillings with Farey determinant $1$.

It fails when $H_1(M_{\mathrm{cusp}}) \cong \ZZ$:
the manifold $\texttt{m004}$ (with $H_1 \cong \ZZ$)
produces a tower whose elements all give cover prime $11$,
independent of position.
\end{theorem}

The essential mechanism: when $H_1(M_{\mathrm{cusp}}) \cong \ZZ/5 \oplus \ZZ$,
the torsion subgroup forces the surgery relation matrix to acquire
determinant $5(ap+bq)$ (for specific peripheral coefficients $a,b$),
and the Farey step $k$ then enters the cover prime formula through
the recursive relator structure.
When $H_1 \cong \ZZ$, the $\ZZ/5$ in the filling comes from the
numerator of the surgery slope, not from the cusped parent,
and the recursive increment mechanism is absent.

%%------------------------------------------------------------------
\section{Conjectures and Open Problems}
\label{sec:conjectures}
%%------------------------------------------------------------------

\begin{conjecture}[Uniform algebraic proof]
\label{conj:induction}
Theorem~\ref{thm:cover} holds for all $k \geq 1$.
The inductive step --- that appending one copy of the chunk
$C = \texttt{abAbaBAB}$ transforms $(\ZZ/\pk)^2$
to $(\ZZ/p_{k+1})^2$ --- follows from showing that the
rank-5 circulant perturbation contributed by $C$ increases
the largest Smith Normal Form invariant by exactly $p_{k+1}-\pk = 10k+6$.
\end{conjecture}

\begin{conjecture}[Universality]
\label{conj:universal}
Theorem~\ref{thm:universality} holds for all cusped hyperbolic
3-manifolds with $H_1 \cong \ZZ/5 \oplus \ZZ$ and Farey
determinant $1$, in particular for the third HFG manifold
$\texttt{m206}$ and its Farey tower.
\end{conjecture}

\begin{conjecture}[Sato-Tate equidistribution]
\label{conj:st}
The Frobenius angles $\theta_k = \arccos(a_{\pk}/(2\sqrt{\pk}))$
for prime $\pk$ are equidistributed in $[0, \pi]$ with respect
to the Sato--Tate measure $(2/\pi)\sin^2\theta\, d\theta$.
\end{conjecture}

The Sato--Tate theorem \cite{CHST2008} guarantees equidistribution
for the full prime sequence; Conjecture~\ref{conj:st} extends this
to the thin subsequence $\{\pk : \pk \text{ prime}\}$.
Numerical evidence through $k = 15$ is consistent with
equidistribution.

\begin{conjecture}[Lucas prime periodicity]
For every Lucas prime $L > 7$, the congruence
$5k(k+1)+1 \equiv 0 \pmod{L}$ has exactly two solutions
modulo $L$, and these two residue classes are symmetric about
$(L-1)/2$.
\end{conjecture}

%%------------------------------------------------------------------
\section{Conclusion}
\label{sec:conclusion}
%%------------------------------------------------------------------

We have established a complete closed chain connecting:
\begin{enumerate}
\item the hyperbolic topology of the PMNS and CKM Farey towers
      (Reidemeister--Schreier, Theorem~\ref{thm:cover});
\item the rational 5-torsion of $\Xo$
      (Proposition~\ref{prop:torsion});
\item the Hecke eigenvalue congruence $a_{\pk} \equiv 2 \pmod{5}$
      (Theorem~\ref{thm:torsion}); and
\item a canonical injection of $\ZZ/5$ into the covering homology
      (Corollary~\ref{cor:shared}).
\end{enumerate}
The $\ZZ/5$ torsion is the organizing principle: it appears in
the base homology of the flavor manifolds, in the rational torsion
of $\Xo$, and in the Hecke eigenvalue sequence, and it forces all
three structures into arithmetic alignment through the single identity
$\pk \equiv 1 \pmod{5}$.

The conductor of $\Xo$ is $11 = p_1$, the first cover prime and
the first Lucas prime that divides any $\pk$.
Both canonical HFG manifolds share this prime as their first
cover prime, providing an arithmetic unification of the lepton
and quark flavor sectors at the conductor scale.

\paragraph{Reproducibility.}
All computations use SnapPy~\cite{SnapPy} and SageMath~\cite{Sage}.
Scripts are available at
\url{https://github.com/drmlgentry/hyperbolic-flavor-scan}.

%%------------------------------------------------------------------
\section*{Acknowledgments}
%%------------------------------------------------------------------

The author thanks the developers of SnapPy and SageMath.

%%------------------------------------------------------------------
\begin{thebibliography}{99}
%%------------------------------------------------------------------

\bibitem{GentrySSRN}
  M.~L.~Gentry,
  Hyperbolic flavor geometry: Standard Model flavor parameters
  from arithmetic hyperbolic 3-manifolds,
  SSRN 6775158, 2026.

\bibitem{GentryEM}
  M.~L.~Gentry,
  Arithmetic invariants of two hyperbolic flavor manifolds:
  WRT invariants, cusped parents, and an $X_0(11)$ bridge,
  submitted to \emph{Experimental Mathematics}, 2026.

\bibitem{GentryFarey}
  M.~L.~Gentry,
  A unified cover prime formula for the PMNS and CKM
  Farey towers of arithmetic hyperbolic 3-manifolds,
  preprint, 2026.

\bibitem{SilvermanAEC}
  J.~H.~Silverman,
  \emph{The Arithmetic of Elliptic Curves}, 2nd ed.,
  Springer, 2009.

\bibitem{CHST2008}
  L.~Clozel, M.~Harris, R.~Taylor,
  Automorphy for some $\ell$-adic lifts of automorphic mod
  $\ell$ Galois representations,
  \emph{Publ. Math. IHES} \textbf{108} (2008), 1--181.

\bibitem{MaR2003}
  C.~Maclachlan, A.~W.~Reid,
  \emph{The Arithmetic of Hyperbolic 3-Manifolds},
  Springer, 2003.

\bibitem{SnapPy}
  M.~Culler, N.~Dunfield, M.~Goerner, J.~Weeks,
  SnapPy, a computer program for studying the geometry and
  topology of 3-manifolds,
  \url{http://snappy.computop.org}.

\bibitem{Sage}
  The Sage Developers,
  SageMath, the Sage Mathematics Software System,
  \url{https://www.sagemath.org}.

\bibitem{NZ1985}
  W.~Neumann, D.~Zagier,
  Volumes of hyperbolic three-manifolds,
  \emph{Topology} \textbf{24} (1985), 307--332.

\bibitem{Diamond2005}
  F.~Diamond, J.~Shurman,
  \emph{A First Course in Modular Forms},
  Springer, 2005.

\end{thebibliography}

\end{document}
